% Ad Atticum 11.14 --- headnote
% ad-atticum-11-14, c. 26 April 47 BC, Brundisium

\textit{Cicero to Atticus, written from Brundisium around
the sixth day before the Kalends of May 47 BC --- about
26 April (the manuscript dateline: \textit{Scr.\ Brundisii
circ.\ vi K.\ Mai.\ a.\ 707 (47)}). Cicero has been
marooned in Brundisium since the previous autumn,
having crossed back to Italy after Pharsalus with
Caesar's permission but no clearance to come further
north or to rejoin public life. Atticus's last letter
has stopped even pretending to console him, and Cicero
does not resent the change: the consolations of the
earlier letters assumed he had company in his fault,
and that assumption is collapsing. The Pompeians who
had taken refuge in Achaia or in Asia --- the men
suing for Caesar's pardon, and even some who had not
been pardoned --- are reported to be sailing for
Africa to join the surviving resistance under Cato
and Scipio. So Cicero will be left with no partner in
his choice except Laelius; and even Laelius, he notes
bitterly, is in a better position, having already been
formally received back.}

\textit{The rest is operational. Cicero assumes Caesar
must have written to Balbus and Oppius about him, and
asks Atticus to sound them out: any deliverance from
that quarter will not be solid, but at least some plan
can be laid. He dreads being seen --- not least with
the son-in-law Dolabella who has come to nothing ---
but in present ills sees no alternative. His brother
Quintus, who has been writing furiously against him,
is reportedly preparing to cross to Africa with the
rest; and Cicero will write to Minucius at Tarentum
about the thirty thousand sesterces he is trying to
raise (the Fufidian estate transactions of the
previous letters are the cover). The closing
sentences are textually broken: two daggered cruxes
are preserved verbatim, and the sense given is the
most natural reading of the corrupted text.}
